%% =====================================================================
%% Jakarta Heat Inequality Mapper -- Zenodo preprint draft
%% Canonical manuscript source. Build with the Multi-Target template:
%%   ./build.sh --target zenodo   --input-dir input --main Manuscript.tex
%%   ./build.sh --target neutral  --input-dir input --main Manuscript.tex
%% Status: WORKING PREPRINT DRAFT -- SCIENTIFIC MODEL DECISIONS PENDING.
%% =====================================================================
\providecommand{\ManuscriptTarget}{neutral}
\documentclass[times,twoside]{_Templates/MultiTargetManuscript}

\usepackage{url}
\usepackage{booktabs}
\usepackage{tabularx}
\usepackage{array}
\usepackage{multirow}
\usepackage{threeparttable}
\usepackage{amsmath}
\usepackage{amssymb}
\usepackage{graphicx}
\usepackage{xcolor}
\usepackage{siunitx}
\usepackage{seqsplit}
\usepackage{float}

% The consolidated class numbers subsections with letters (A, B, ...).
% Use conventional decimal numbering for this manuscript without changing
% the shared class or its target profiles.
\renewcommand{\thesubsection}{\thesection.\arabic{subsection}}
\renewcommand{\thesubsubsection}{\thesubsection.\arabic{subsubsection}}

% The consolidated class is a two-column bioRxiv-style layout and labels
% numbered sections as "Supplementary Note". This manuscript is a one-column
% research article with conventional section numbering. The following
% manuscript-level overrides do not modify the shared class or its profiles.
\titleformat{\section}[hang]{\large\sffamily\bfseries}{\thesection.}{0.5em}{#1}[]
\titleformat{\subsection}[hang]{\normalsize\sffamily\bfseries}{\thesubsection.}{0.5em}{#1}[]
\titleformat{\subsubsection}[hang]{\small\sffamily\bfseries\itshape}{\thesubsubsection.}{0.5em}{#1}[]
\titlespacing*{\section}{0pc}{2.6ex plus 4pt minus 3pt}{1.1ex}
\titlespacing*{\subsection}{0pc}{2.1ex plus 3pt minus 2pt}{0.7ex}
\titlespacing*{\subsubsection}{0pc}{1.8ex plus 2pt minus 1pt}{0.6ex}
\makeatletter
\@tmptwocolumnfalse
\makeatother

% Readability adjustments for a one-column research article. The class's
% abstract is bold by default; normal weight is easier to read at this length.
\setlength{\parindent}{1.2em}
\setlength{\parskip}{0.25em}
\renewcommand{\abstract}{\small\normalfont}
\setlength{\bibsep}{1.5pt}
\setlength{\belowcaptionskip}{4pt}

\leadauthor{Faqih}

\begin{document}
\onecolumn
\raggedbottom
\setlength{\emergencystretch}{2.5em}

\title{From Thermal Hotspots to Spatial Priorities: Attribution and Sensitivity of Urban Heat Screening in Jakarta}
\shorttitle{Attribution and sensitivity of Jakarta heat priorities}

\hypersetup{
  pdftitle={From Thermal Hotspots to Spatial Priorities: Attribution and Sensitivity of Urban Heat Screening in Jakarta},
  pdfauthor={Muhammad Fariq Faqih; Muhammad Alim Mustofa; Mu'adz Abdulloh},
  pdfsubject={Working preprint draft -- scientific model decisions pending},
  pdfkeywords={urban heat; land surface temperature; environmental inequality; park accessibility; composite index; sensitivity analysis; spatial autocorrelation; Jakarta}
}

% Author identities and affiliations supplied by the project coordinator.
% ORCIDs remain to be added if the authors wish to include them.
\author[1]{Muhammad Fariq Faqih}
\author[1]{Muhammad Alim Mustofa}
\author[1]{Mu'adz Abdulloh}
\affil[1]{Renaro Lab, student-led independent AI laboratory, \url{https://renaro.my.id}}

\maketitle

\begin{center}
\fbox{\parbox{0.94\textwidth}{\centering\small
\textbf{WORKING PREPRINT DRAFT --- SCIENTIFIC MODEL DECISIONS PENDING.}
This manuscript is an internal scientific draft prepared for Zenodo-style
preprint review. It is not a peer-reviewed article, not an approved scientific
release, and not a policy recommendation. Model weights, thresholds, sample
eligibility, access semantics and spatial-inference conventions remain
provisional and unsigned. Numerical results are traceable to the frozen RC1
artifact described in the Data and Code Availability section. The Zenodo profile
in the footer identifies the intended presentation format only; it does not
indicate a deposit, submission, acceptance or DOI.
}}
\end{center}

\begin{abstract}
\noindent\textbf{Abstract.}\quad Urban heat mitigation planning increasingly relies on composite spatial
indices that combine land surface temperature (LST), vegetation, population and
access to green space. Such indices produce priority maps, but the maps are
model outputs rather than observed inequalities, and their dependence on
demographic and accessibility assumptions is rarely quantified. We examine how
spatial priorities change across four screening models in the five mainland
administrative cities of Jakarta, Indonesia, on a 500~m grid (2,765 cells).
M1 ranks LST only; M2 adds vegetation; M3 adds modelled total population; and
M4 adds a count of children and older adults plus a modelled walking-network
distance to 1,015 OpenStreetMap candidate parks. The common comparison sample
contains 2,536 cells (91.7\% of the master grid). Thermal and demographic inputs
refer to a declared 2024 period; the park and pedestrian-network snapshot is from
19 November 2026, so the analysis is a mixed-period candidate proxy. The M1 and M4 top-decile sets share 141 of 254 cells (55.5\% retention), but
their Jaccard overlap is 0.384 because 113 cells are unique to one of the two
sets; the mean absolute M1-to-M4 rank shift is 307 positions. Between M3 and M4, an
exact three-factor Shapley decomposition attributes 85.5\% of absolute rank
movement to park addition, 13.2\% to reweighting and 1.3\% to the age-count
substitution. The age count is nearly collinear with total population
(Spearman $\rho=0.9995$); replacing it with total population leaves the top
decile almost unchanged (Jaccard 0.977), whereas replacing it with an age share
changes the set substantially (0.347). Under the provisional additive M4, the
top-decile cutoff (0.8737) exceeds the maximum score of all 537 common cells
within 500~m of a candidate park (0.8696); no near-park cell enters the top 254,
the closest top cell is 638~m from a candidate park, and 174 of 254 top cells are
at the 1,000~m distance cap. This is a property of the index and its observed
inputs, not evidence that parks are unimportant. Access assumptions matter:
a positive-tag-only destination stress test yields Jaccard 0.420 with the
baseline top decile, while a conservative connector scenario yields 0.896.
The 229 excluded cells are cooler, greener and more peripheral than included
cells, and the common sample is therefore a conditional estimand rather than a
Jakarta-wide estimate. Local Moran labels are convention-sensitive: 663 cells
are significant under the provisional directed/999/Benjamini--Hochberg
convention, 496 under two-sided testing, and zero under Benjamini--Yekutieli at
999 permutations, where the minimum attainable permutation $p$-value exceeds the
correction threshold. We make no claim about health outcomes, verified public
park access, historical 2024 network conditions or intervention effectiveness.
The contribution is an attribution and uncertainty audit of a reproducible
screening index, including negative findings about age-count redundancy and the
mechanical behaviour of an additive park-deficit term.
\end{abstract}

\begin{keywords}
\textbf{Keywords:}\quad urban heat \textbar{} land surface temperature \textbar{} environmental inequality \textbar{} park accessibility \textbar{} composite index \textbar{} sensitivity analysis \textbar{} spatial autocorrelation \textbar{} Jakarta
\end{keywords}

\begin{corrauthor}
fariq\at renaro.my.id;\ alim\at renaro.my.id;\ muadz\at renaro.my.id
\end{corrauthor}

\section{Introduction}

Urban heat is a planning problem before it is a measurement problem. Cities
must decide where to plant trees, open cooling spaces, retrofit buildings or
target public health outreach, and these decisions require some ordering of
places. Satellite land surface temperature (LST) is attractive for that task
because it is spatially continuous, repeatable and available at metropolitan
scale \citep{oke1982,voogt2003}. LST has therefore become a common first layer in
urban heat assessments, including in Jakarta \citep{siswanto2023,putra2021,rizki2024}.
Yet LST is a surface property retrieved from clear-sky remote sensing; it is not
air temperature, personal heat exposure or a health outcome
\citep{schwarz2012,voogt2003}. A map of the hottest pixels is a map of one
hazard component, not a map of where people are most exposed or where an
intervention would do the most good.

Planning questions usually require more than a hazard layer. Exposure asks how
many people are present; vulnerability asks whether some groups face greater
risk for a given exposure; and accessibility asks whether people can reach
cooling or green space \citep{reid2009,wolf2013,harlan2006}. These are distinct
constructs. Combining them into a single composite index is a legitimate
decision-support strategy, but the combination is not a neutral act: it encodes
a theory of what matters, a set of weights, a normalisation rule and a
definition of access \citep{malczewski2006,saisana2005,greco2019}. Two indices
can use the same data and produce different priority maps because they answer
different questions. The policy relevance of a composite therefore depends on
how transparently those choices are exposed and how strongly the resulting
ranking depends on them.

Jakarta provides a demanding setting for this problem. The metropolitan region
is hot, humid and highly unequal in built form and green-space provision. LST
and vegetation studies document strong intra-urban thermal contrasts
\citep{siswanto2023,putra2021,rizki2024}, while demographic and epidemiological
work shows that older adults face elevated heat-related mortality risk
\citep{varquez2020}. Heat is also experienced through social practice and
everyday access to shade, water and communal space, particularly in informal
settlements \citep{salsabila2023}. Park access is uneven and cannot be inferred
from park counts alone; it depends on the pedestrian network, entrances,
permissions and the quality of destinations \citep{wolch2014,rigolon2016,ekkel2017}.
Jakarta therefore offers a case where a heat-only map is clearly incomplete, but
where the information needed to complete it is itself uncertain.

A growing literature combines heat, demographic vulnerability and green-space
access. Sidiqui et al.\ \citep{sidiqui2022} map urban heat vulnerability with
GIS data and socio-demographic indicators; Kim and Kwak \citep{kim2026} develop a
cooling-weighted park-accessibility framework for older adults using an
OpenStreetMap walking network; Fung et al.\ \citep{fung2024} prioritise social
vulnerability in urban heat mitigation; and Sobhaninia et al.\
\citep{sobhaninia2025} review methods for siting green infrastructure. The idea
of combining LST, age and access is therefore not new. What remains
underdeveloped is a quantitative account of how much each added construct
changes a priority ranking, how much of the change is attributable to
accessibility rather than reweighting, and how sensitive the result is to
imperfect destination and network data. Those questions are methodological and
empirical at once, and they are the focus of this study.

We ask four questions. First, how much do priority rankings and top-decile sets
change when LST is supplemented with vegetation, population, age count and
modelled park distance? Second, what share of the M3-to-M4 change is attributable
to park addition, reweighting and the age substitution? Third, how sensitive are
the priorities to plausible assumptions about park destinations, connector
permissions, missing routes and demographic specification? Fourth, how do sample
eligibility and local spatial-inference conventions limit what can be inferred
from the resulting map? We answer these questions using the frozen RC1 artifact
of the Jakarta Heat Inequality Mapper, a reproducible 500~m grid analysis of the
five mainland cities \citep{jhim2026}.

The paper makes four contributions. Methodologically, it provides an exact
attribution and uncertainty audit of a composite heat-screening index, including
a rank-based Shapley decomposition and a scenario envelope that separates
construct changes from reweighting. Empirically, it shows for Jakarta that an
age-count term adds little information beyond total population, while an
age-share term changes the objective and the geography; and that an additive
park-deficit term can mechanically determine the top decile under observed
inputs. Critically, it documents a negative result: a park term can dominate a
ranking without demonstrating that parks are unimportant or that public access
is unequal. Finally, it provides a reproducible workflow and a claim-evidence
register that distinguish verified arithmetic from provisional interpretation.
We do not claim that M4 is a better model than M1--M3, that the top-priority
cells are the most heat-vulnerable places in Jakarta, or that the results
demonstrate the health benefits of any intervention.

Section~2 situates the work in prior literature. Section~3 describes the study
area and data. Section~4 gives the complete methodology. Section~5 reports
results, Section~6 discusses their interpretation, and Section~7 states
limitations. Section~8 concludes. Supplementary material contains additional
sensitivity tables, a data-quality register and figure provenance.

\section{Related Work}

\subsection{Land surface temperature and the urban heat island}

The urban heat island is a well-established phenomenon in which built surfaces
and reduced vegetation alter the surface energy balance \citep{oke1982}. Thermal
remote sensing makes the surface expression of that balance observable at fine
spatial resolution \citep{voogt2003}. LST is strongly related to vegetation
abundance: greener and more vegetated surfaces tend to be cooler, and the
LST--NDVI relationship has been used widely to map surface urban heat islands
\citep{weng2004,peng2012,estoque2017}. Satellite platforms also reveal that
surface urban heat island intensity varies substantially across cities and
regions \citep{peng2012,zhou2019}. These studies justify LST and NDVI as
screening layers. They also caution against equating LST with air temperature or
human heat stress; the relationship between surface and air temperature is
context-dependent \citep{schwarz2012}. Our M1 and M2 models are deliberately
simple representations of this tradition.

\subsection{Exposure, vulnerability and composite indices}

Heat vulnerability indices combine hazard with demographic sensitivity and
adaptive capacity. Reid et al.\ \citep{reid2009} mapped community determinants of
heat vulnerability in the United States, and Wolf and McGregor
\citep{wolf2013} developed a heat-wave vulnerability index for London. Such
indices are useful for communication and screening, but they are sensitive to
indicator selection, weighting, normalisation and aggregation
\citep{saisana2005,paruolo2013,greco2019}. A composite can change rankings even
when the underlying data are unchanged. This is not a defect unique to heat
indices; it is a general property of multidimensional indicators.

A parallel literature examines inequality in heat exposure. Chakraborty et al.\
\citep{chakraborty2019} found disproportionately higher exposure to urban heat
in lower-income neighbourhoods across multiple cities; Hsu et al.\
\citep{hsu2021} documented unequal exposure to urban heat island intensity across
major US cities; and Tuholske et al.\ \citep{tuholske2021} quantified global
urban population exposure to extreme heat. Heat-health risk frameworks combine
hazard, exposure and social vulnerability in related ways \citep{dong2020}.
These studies establish that exposure is uneven, but they generally do not model
walking access to cooling destinations or decompose the effect of an access term
on a priority ranking. Our work is closer to a screening and sensitivity
exercise than to an epidemiological exposure assessment.

\subsection{Age, heat and green-space access}

Children and older adults are often treated as heat-vulnerable groups. Reviews
of temperature and mortality report elevated risk among older adults
\citep{basu2009,astrom2011}, and children's outdoor microclimates raise distinct
physiological and behavioural concerns \citep{vanos2015}. Global projections
suggest that heat exposure of older adults will increase substantially
\citep{falchetta2024}. In Jakarta, Varquez et al.\ \citep{varquez2020} projected
future increases in elderly heat-related mortality under urbanisation and
climate change. These findings motivate an age-sensitive screening layer. They
do not imply that an age \emph{count} is a measure of vulnerability composition:
counts largely track population size, whereas shares describe the relative
presence of vulnerable groups.

Green-space access is similarly multidimensional. Reviews document inequities
in park provision and quality \citep{wolch2014,rigolon2016}, and accessibility
metrics differ in whether they measure distance, travel time, network reach or
perceived access \citep{ekkel2017}. GIS studies have long used network distance
to assess park distribution \citep{oh2007}; spatial-accessibility methods provide
a broader family of distance, catchment and floating-catchment measures
\citep{luo2003,luo2009}, and walking-network studies make the route itself the
object of analysis \citep{stoia2022}. More recent work incorporates cooling
weights and older-adult demand \citep{kim2026}. Feng et al.\
\citep{feng2019} analyse the distribution and accessibility of urban parks in
Beijing, and Zhang et al.\ \citep{zhang2024} examine park size thresholds,
cooling accessibility and equity. Park cooling itself depends on size,
vegetation, configuration and distance \citep{bowler2010,aram2019,masoudi2019,algretawee2022}.
Access to a cooling refuge is also an equity question in its own right
\citep{voelkel2018,amorim2023}. The broader point is that a
park-distance term is a model of access, not a measurement of public access.
OSM completeness and tagging quality vary, and network routing depends on
connectors and permissions that are rarely verified \citep{haklay2010,barron2014,herfort2023}.
Our treatment of candidate parks as a proxy, and our explicit stress tests, follow
from this literature.

\subsection{Sensitivity, robustness and spatial inference}

Composite-indicator research has developed a rich toolkit for uncertainty and
sensitivity analysis \citep{saisana2005,paruolo2013,greco2019}. The central
lesson is that rankings should be reported with the assumptions that generate
them, not as unconditional facts. In parallel, spatial-analytic work provides
global and local measures of autocorrelation \citep{anselin1995,getis1992}, but
local inference is sensitive to the choice of weights, tail specification,
permutation count and multiple-testing correction
\citep{benjamini1995,benjamini2001}. A cluster map is therefore not a
self-interpreting result. Our design treats model comparison, attribution,
scenario analysis and LISA sensitivity as parts of one evidentiary system rather
than as separate appendices.

\subsection{Jakarta-specific evidence}

Jakarta has a substantial applied literature on LST, vegetation and heat. Putra
et al.\ \citep{putra2021} document increasing urban heat island area in relation
to land-use change; Siswanto et al.\ \citep{siswanto2023} characterise the
spatio-temporal urban heat island of the Jakarta metropolitan area; and Rizki et
al.\ \citep{rizki2024} examine the impact of urban green space on the heat island
in East Jakarta. Salsabila et al.\ \citep{salsabila2023} show through ethnographic
and practice-based work that cooling in Jakarta's informal settlements is
entangled with social life and communal space. Manik and Syaukat
\citep{manik2017} compare heat vulnerability in two tropical Indonesian cities.
These studies establish the local relevance of heat and vegetation, but they do
not provide a cell-level attribution of how demographic and network-access
assumptions reshape a priority ranking. That is the space this paper occupies.

\subsection{Positioning and contribution}

Table~\ref{tab:literature} summarises the closest bodies of work. The table is
not a systematic review and does not establish that no similar study exists. It
shows that every individual ingredient of our analysis has precedent: LST--NDVI
mapping, heat vulnerability indices, age-sensitive risk, network-based park
access and composite-indicator sensitivity. The combination-only claim is
therefore not defensible. What remains is a narrower and more useful claim: a
quantitative attribution and uncertainty audit of how priorities shift when
these ingredients are assembled in one Jakarta screening index, including
negative findings about redundancy and mechanical threshold behaviour. We
present that claim as provisional and conditional on the frozen RC1
configuration.

\input{Tables/table2_literature}

\section{Study Area and Data}

\subsection{Study area and spatial unit}

The study covers the five mainland administrative cities of Jakarta, Indonesia:
West Jakarta (\textit{Jakarta Barat}), Central Jakarta (\textit{Jakarta Pusat}),
South Jakarta (\textit{Jakarta Selatan}), East Jakarta (\textit{Jakarta Timur})
and North Jakarta (\textit{Jakarta Utara}). Kepulauan Seribu, the offshore
administrative regency, is excluded. The analysis unit is a 500~m lattice in the
projected UTM zone 48S coordinate reference system (EPSG:32748), clipped to the
five-city land boundary. The master grid contains 2,765 cells. Because the
lattice is clipped to a complex coastline and administrative boundary, 383 cells
are partial polygons; 65 have less than 10\% of the nominal 0.25~km$^2$ area, and
the smallest is approximately $4.2\times10^{-7}$~km$^2$. These slivers are
retained and flagged rather than silently deleted, because deleting them would
change the population accounting and the spatial graph.

The grid is a spatial sampling device, not a cadastral or neighbourhood unit.
Cell-level results are therefore subject to the modifiable areal unit problem:
different partitions of the same land could produce different rankings and
different local statistics \citep{openshaw1984}. We do not treat the 500~m cell
as a meaningful social unit. It is a common, reproducible support on which
heterogeneous raster and vector inputs can be compared.

\subsection{Thermal and vegetation data}

LST and NDVI were derived from Landsat 9 Collection 2 Level-2 surface
reflectance and surface temperature products (dataset
\texttt{LANDSAT/LC09/C02/T1\_L2}) \citep{usgs2020} for a declared 2024 window
(1 January to 31 December 2024). The intended recipe filters to L2SP products, masks fill,
dilated cloud, cirrus, cloud, cloud shadow and snow using QA\_PIXEL bits 0--5,
requires QA\_RADSAT = 0, converts ST\_B10 to Celsius using the Collection 2
scale factors, computes NDVI from SR\_B4 (red) and SR\_B5 (near-infrared), and
takes a per-pixel median composite across clear observations. Each band is
reduced independently, so the LST and NDVI medians for a pixel may derive from
different acquisition dates. The composite is exported at 30~m in EPSG:32748.

Dynamic World class probabilities are available for 2,764 cells
\citep{brown2022}, but categorical class fractions were not exported; M1--M4 do
not use them, and we make no categorical land-cover claim.
Grid values are area-weighted means over valid pixels within each clipped cell,
computed with exactextract. A cell is flagged for low thermal support when the
valid LST area fraction is below 0.50; this threshold is provisional. One cell
(\texttt{JKT\_R037\_C062}) has 27.759\% valid LST area (6,724.1~m$^2$ of
24,223.1~m$^2$) and is excluded from all models under the 0.50 rule. The
delivered LST and NDVI values for that cell are retained in the source table, but
the cell does not enter the common comparison sample. We do not interpret
low-support cells as having no heat.

The critical provenance limitation is that the original Earth Engine scene,
task and AOI records could not be recovered. The scene manifest is header-only;
this means that records are unavailable, not that zero scenes were acquired.
The 2024 label is therefore a \emph{declared recipe period}, not an
independently verified acquisition history. Annual clear-observation counts are
available, but monthly and seasonal sampling cannot be reconstructed from the
median raster alone. We make no claim of seasonal representativeness, annual mean
temperature, temporal trend or a verified 2024 snapshot.

\subsection{Population and age data}

Population inputs are from the WorldPop Global2 R2025A release, which provides
modelled people per grid cell at approximately 100~m (3 arc-second) resolution
for a 2024 reference year, including total population and age/sex groupings
\citep{worldpop2025}. We use total population, children aged 0--14 and older
adults aged 65 and over. These are modelled estimates, not a census or a
register. They describe residential population, not daytime presence, mobility
or personal exposure.

Zonal aggregation uses the project's hybrid rule: where more than one valid
source-pixel centre falls inside a clipped cell, the valid centre values are
summed; otherwise, intersecting valid pixels are allocated fractionally by
their intersection area. A fractional alternative allocates every valid native
pixel by its polygon intersection fraction and is additive under changes of
partition. The hybrid rule is not additive under repartitioning because centres
may be selected differently in each polygon. The primary analysis uses the
hybrid rule for continuity with the frozen baseline; the fractional alternative
is a sensitivity scenario. Population counts are never multiplied by coverage a
second time.

Source support is treated explicitly. Of 19 cells that originally had no valid
age-raster support, 13 were recovered by retrieving buffered age surfaces that
cover the five-city union with a 1,000~m buffer. Six cells remain unsupported in
total population and both age groups:
\texttt{JKT\_R003\_C031}, \texttt{JKT\_R057\_C063},
\texttt{JKT\_R061\_C009}, \texttt{JKT\_R061\_C010},
\texttt{JKT\_R061\_C011} and \texttt{JKT\_R062\_C043}. They are retained in the
master grid with analytical nulls, while their original exported zeros are
preserved in \texttt{*\_raw} columns. A null is not a estimated population of
zero. Two of the six unsupported cells are substantively large
(\texttt{JKT\_R061\_C009}: 0.244~km$^2$, 97.8\% coverage;
\texttt{JKT\_R061\_C010}: 0.099~km$^2$, 39.5\% coverage), both in North Jakarta.

\subsection{Candidate parks and the pedestrian network}

Park destinations are derived from an OpenStreetMap snapshot extracted on
19 November 2026. The inventory contains 1,015 candidate park polygons. Each
candidate was reviewed against archived OSM tags and, where available, official
or operator context. No candidate has verified public access and no candidate has
a verified physical entrance: 980 are classified as unknown, 28 as conditional
and unverified, and 7 as carrying a positive OSM access tag that remains
unverified. All 1,015 are unresolved in the evidence register. A positive OSM
tag is an assertion, not a field verification; an unknown tag is not evidence of
closure. The candidate set is therefore a destination proxy, not a public-park
inventory. Table~\ref{tab:park} summarises the evidence states.

Walking distances were computed on a pedestrian graph built from the same OSM
snapshot. Each grid cell and each candidate park is connected to the graph at a
nearest node; the park is represented by a representative point inside its
polygon, not by a verified entrance. The baseline distance is the shortest-path
network distance to any non-restricted candidate park, capped at 1,000~m for the
park-deficit term. The cap is a modelling choice, not a measured threshold of
park utility. Of 2,765 master cells, 226 have no route to a candidate park in
the baseline graph. These are retained as missing and excluded from the common
comparison sample; they are not assigned a distance of zero, infinity or the cap.

Two route-permission stress scenarios are used. A conservative scenario removes
edges incident on, or matching, explicitly forbidden source nodes
(\texttt{foot=no/private}, or \texttt{access=no/private} without an explicit
permissive foot override) and excludes direct connectors that touch those nodes.
It removes 24,370 directed edges and increases missing routes from 226 to 407.
A positive-tag-only destination scenario retains only destinations with a
positive access or permissive tag; it contains 10 destinations (9 tagged
\texttt{yes} and 1 \texttt{permissive}), of which 8 are within 500~m and 41
within 1,000~m of a routable origin; 2,482 cells are saturated at the 1,000~m
cap. This scenario is a coverage stress test, not a better estimate of access.
In the connector audit, 167 origin connectors and 78 park connectors conflict
with explicit OSM node restrictions. These conflicts are retained in the
evidence file and are not silently repaired.

\input{Tables/table10_park}

\subsection{Reference periods and source table}

The analysis mixes reference periods. LST and NDVI refer to the declared 2024
thermal period; WorldPop refers to a 2024 reference year; and the park and
pedestrian network is a 19 November 2026 snapshot. We do not have a historical
OSM archive: six requests for historical snapshots failed, and no temporal
comparison is possible. The resulting object is a \emph{mixed-period candidate
proxy}. It can support statements about how a 2024 thermal/demographic pattern
combines with a 2026 candidate-access network under explicit assumptions. It
cannot support statements about 2024 public access, seasonal heat exposure,
temporal change or the effects of a 2026 network on 2024 residents.
Table~\ref{tab:sources} summarises the sources, their processing and their main
limitations.

\input{Tables/table1_sources}

\section{Methods}

\subsection{Sample eligibility and the common comparison set}

All models are compared on a common set of cells for which every required input
is available. This avoids the common error of comparing correlations or top-set
overlap across different sample sizes. Starting from 2,765 master cells, 229 are
excluded: 222 have a missing walking route only; 4 have missing route,
population and age support; 2 have missing total population and age support but a
known route; and 1 has low LST support under the provisional 0.50 rule. The
union is 229, because the route-missing and support-missing categories overlap.
The common set is therefore 2,536 cells (91.7\%; Table~\ref{tab:eligibility}).
Model-specific eligibility is 2,764 for M1 and M2, 2,758 for M3 and 2,536 for M4. M1 and M2 differ from the
master grid only by the low-LST cell; M3 additionally loses six cells without
population support; M4 additionally loses the route-missing and age-unsupported
cells.

The common set is not a random sample of Jakarta. The excluded cells are
systematically cooler, greener, more peripheral and less connected than the
included cells. We report their profile in Section~5.8 and treat the common-set
result as a conditional estimand. For spatial statistics, one additional cell
(\texttt{JKT\_R013\_C019}) is a Queen-contiguity island; it is removed from the
graph, leaving 2,535 observations and 18,368 directed edges. This island is a
graph property, not a source-support exclusion.

\input{Tables/table4_eligibility}

\subsection{Zonal aggregation and normalisation}

For a cell $i$ and a variable $x$, let $x_i$ be the area-weighted mean over valid
source pixels, or the modelled population count allocated to the cell. Each
model component is normalised to $[0,1]$ using a percentile-clipped min--max
rule computed on the common sample:
\begin{equation}
n(x_i) = \operatorname{clip}\!\left(
\frac{x_i - q_{0.05}(x)}{q_{0.95}(x) - q_{0.05}(x)},\ 0,\ 1
\right),
\label{eq:norm}
\end{equation}
where $q_{0.05}$ and $q_{0.95}$ are the 5th and 95th percentiles of the
component in the common sample. Values below the 5th percentile are set to zero
and values above the 95th percentile to one. The rule reduces the influence of
extreme cells and makes components comparable, but it also saturates scores:
M1 has 127 cells tied at the maximum. Ties are broken deterministically by the
unclipped min--max score and then by grid identifier, so the ranking does not
depend on row order. The park-deficit component is not renormalised after
capping; it is defined directly on a 0--1 scale as
$\min(d_i,1000)/1000$, where $d_i$ is the modelled network distance in metres.

\subsection{The four screening models}

The four models are nested in data requirements but not in statistical
estimation. They are alternative constructs, not successive estimates of one
latent quantity. Let $L_i$ be LST, $V_i$ NDVI, $P_i$ total population,
$C_i$ the population aged 0--14, $E_i$ the population aged 65 and over, and
$D_i$ the candidate-park network distance. The scores are
\begin{align}
S_{1,i} &= n(L_i), \label{eq:m1}\\
S_{2,i} &= 0.60\,n(L_i) + 0.40\,[1-n(V_i)], \label{eq:m2}\\
S_{3,i} &= 0.40\,n(L_i) + 0.30\,[1-n(V_i)]
          + 0.30\,n(\log(1+P_i)), \label{eq:m3}\\
S_{4,i} &= 0.30\,n(L_i) + 0.20\,[1-n(V_i)]
          + 0.25\,n(\log(1+C_i+E_i))
          + 0.25\,\frac{\min(D_i,1000)}{1000}. \label{eq:m4}
\end{align}
M1 is a heat-only screening model. M2 adds a vegetation-deficit term. M3 adds
total population exposure. M4 replaces total population with a count of children
and older adults and adds the park-deficit term. The weights are provisional:
they were fixed before the real LST data were available, are not derived from a
Jakarta-specific literature, and have not been approved by the research
coordinator. They are therefore reported as a configuration, not as a validated
standard. The model formulas are reproduced exactly from the frozen
\texttt{model\_specification.json}; M4 is not silently redefined in this paper.

Table~\ref{tab:models} summarises the four score definitions.
M4 is best described as a \emph{provisional heat--population--park-deficit
screening index}. It is not an age-vulnerability model and not a health-risk
model. The age term is a count, not a share. Because children and older adults
are a nearly constant fraction of the modelled population in this dataset, the
count term is almost collinear with total population. A separate age-share
scenario is used to show what a composition-oriented question would look like;
it is not the primary M4 specification.

\input{Tables/table3_models}

\subsection{Ranking, top-priority sets and model comparison}

Cells are ranked in descending score, with rank 1 assigned to the highest score.
The top-priority budget is $\lceil 0.10 \times 2{,}536\rceil = 254$ cells, i.e.\
the top decile. We also report 5\% (127 cells) and 20\% (508 cells) budgets to
show threshold dependence. Two models are compared using Spearman's rank
correlation, Kendall's $\tau_b$, the Jaccard overlap of their top-$k$ sets,
and the mean and median absolute rank shift. For top sets $A$ and $B$,
\begin{equation}
J(A,B) = \frac{|A \cap B|}{|A \cup B|}.
\label{eq:jaccard}
\end{equation}
Jaccard is 0 for disjoint sets and 1 for identical sets. It is sensitive to the
budget: a small top set can have low overlap even when the underlying scores are
highly correlated. We therefore report Jaccard at three budgets
(Table~\ref{tab:comparison}) and never treat a single value as a general measure
of model agreement.

\subsection{Attribution of M3-to-M4 change}

The M3-to-M4 difference combines three changes: the addition of a park-deficit
term, the reweighting of the shared LST and NDVI terms, and the substitution of
the age count for total population. Let the score change be
\begin{equation}
\begin{split}
\Delta S_i &= S_{4,i} - S_{3,i} \\
&= \underbrace{\left(-0.10\,n(L_i)-0.10\,[1-n(V_i)]-0.05\,n(\log(1+P_i))\right)}_{\text{reweighting}} \\
&\quad + \underbrace{0.25\left[n(\log(1+C_i+E_i))-n(\log(1+P_i))\right]}_{\text{age substitution}}
+ \underbrace{0.25\,\frac{\min(D_i,1000)}{1000}}_{\text{park addition}}.
\end{split}
\label{eq:delta}
\end{equation}
The three terms sum exactly to the score difference. To attribute rank movement
rather than score movement, we compute the exact three-factor Shapley value of
the induced ranking over all $2^3$ coalitions. For each cell, the marginal rank
gain of adding a factor is evaluated in every order, and the average is taken.
The absolute Shapley contributions are then normalised to sum to 100\%. This is
a descriptive decomposition of a model change, not a causal effect of parks,
age or weights on real-world outcomes.

We also run matched-weight ablations that hold the sample and normalisation
anchors fixed. These include M3, a reweighting-only variant, an age-swap-only
variant, a park-addition-only variant, all pairwise combinations, the full M4,
a no-park variant, a park-only variant, a heat-only variant, an equal-weight
variant, a total-count swap, a vulnerable-share variant, child-share and
elderly-share variants, and a proportional M3-plus-park variant. These
ablations answer different questions and are not all equally plausible policy
alternatives; Tables~\ref{tab:attribution_m3} and~\ref{tab:attribution_m4}
report the main results and state their baselines explicitly.

\subsection{Park-access and demographic sensitivity scenarios}

Park-access uncertainty is evaluated with four families. First, destination
scenarios vary the candidate set: baseline, positive-tag-only, and alternative
nearest-distinct-park assignments. Second, permission scenarios vary the graph:
baseline, conservative node/connector removal, and a positive-tag stress
variant. Third, missing-route scenarios assign no distance (baseline exclusion),
the observed median distance, or the 1,000~m cap; these bound the effect of
unknown routes but are not estimates of true access. Fourth, distance-cap
scenarios vary the cap (750, 1,000 and 1,500~m) and the park threshold (500 and
1,000~m). For each scenario we recompute normalisation on the scenario's
eligible set and report Jaccard against the baseline M4 top decile. We
distinguish an \emph{eligibility-changing} comparison, in which the sample
changes, from a \emph{fixed-anchor} comparison, in which the same cells and
normalisation anchors are used. The latter isolates rank changes from sample
changes. The resulting family envelope is reported in Section~5.10.

Demographic sensitivity compares count and composition specifications. The
total-count swap replaces the age count with total population at the same 0.25
weight. The vulnerable-share, child-share and elderly-share variants replace the
count with a share, with the same weight and a positive-population denominator.
These are different estimands: a count asks where many vulnerable people live; a
share asks where vulnerable people are a larger fraction of residents. The
primary M4 is a count-based exposure index, and we do not describe it as an
age-composition model.

\subsection{Global and local spatial autocorrelation}

Global Moran's $I$ is computed on the common sample using Queen contiguity,
row-standardised weights, after removing islands. The null is spatial
randomisation of values; the alternative is positive autocorrelation. The
permutation $p$-value is
\begin{equation}
p_{\mathrm{sim}} = \frac{1 + \#\{I_{\mathrm{perm}} \ge I_{\mathrm{obs}}\}}{R+1},
\label{eq:moranp}
\end{equation}
with $R=999$ and seed 42 in the baseline. This is an upper-tail convention.
Normal-theory $p$-values are reported separately and are not substituted for the
permutation result. At $R=999$, the minimum attainable $p$-value is 0.001; a
reported 0.001 is a resolution limit, not evidence of $p=0$.

Local indicators of spatial association (LISA) use conditional randomisation
with the same weights and seed. The provisional baseline uses the directed
alternative, 999 permutations and Benjamini--Hochberg (BH) correction at
$q=0.05$ across all non-island cells in the M4 map. We report four alternatives:
directed versus two-sided tails, 999 versus 9,999 permutations, and BH versus
Benjamini--Yekutieli (BY) correction. BH controls the false discovery rate under
independence or positive dependence; BY is more conservative under arbitrary
dependence. The testing family is each preselected map separately, not a family
across all interactive selections. A local cluster label is not evidence that
the index is valid, that the cluster is a social process, or that an
intervention would be effective. Table~\ref{tab:lisa} reports the convention
range.

\subsection{Scenario envelope and reproducibility}

The scenario registry contains 52 scenarios in families covering weights
(9 scenarios), park thresholds (4), normalisation (3), QA exclusions (3), route
policy (5), thermal support (2), land cover (1), population specification (2),
joint factorial combinations (16) and other variants. Two families
(entrance/barrier and thermal sampling) are blocked by missing evidence; one
(temporal) is conditional. We report the minimum and maximum Jaccard within each
family at the 10\% budget. The joint factorial is a 16-scenario subset crossing
construct, cap, normalisation and aggregation choices; route-policy variants are
a separate family. We therefore do not claim all-family robustness. The purpose
of the envelope is to show which conclusions survive plausible alternatives and
which do not.

The analysis is reproduced from the frozen RC1 artifact at commit
\texttt{\seqsplit{3b8ed05be92c35bbb2b3a617688540ee856b11c8}}. The implementation SHA-256
over 80 declared files is
\texttt{\seqsplit{6b2eae9dc37b14d92bd6accafda602da1e42af643ba8a72b6431d60e985b233e}}.
The artifact includes 2,765 master grid rows, 2,536 common eligible rows, 254
top-decile rows per model, 63 replayed tables with maximum numeric difference
0.0, and a full test suite of 236 tests plus 7 subtests in the integration
workspace. RC1's original clean-checkout audit found one portability failure
because a raw WorldPop raster is not tracked. The RC2 engineering candidate
remediates that failure: a clean checkout reports 236 tests plus 7 subtests with
0 failures, 0 errors and 0 skips; RC1 and RC2 receipts each validate 9/9 under a
canonical Git-blob SHA-256 policy; the standard and remediation output manifests
validate 69/69 and 80/80 relative paths; and the 63-table replay again reports a
maximum numeric difference of 0.0. RC2's model definitions are byte-identical to
RC1 and remain provisional and unsigned. RC2 is an engineering remediation, not
a new scientific model; all empirical results in this paper remain tied to RC1.
The reproduction commands, input hashes and output hashes are documented in the
artifact handoffs and are summarized in the Data and Code Availability section.

\section{Results}

\subsection{Environmental, demographic and accessibility distributions}

The common sample has a mean LST of 41.49~$^\circ$C (SD 2.29; range
31.28--47.92), a mean NDVI of 0.315 (SD 0.106; range $-0.024$--0.795), and a
modelled population of 10.52 million people (mean 4,148 per cell; median 4,288).
Children and older adults account for 3.11 million people, or 29.5\% of the
modelled population, with a between-cell share that is remarkably narrow
(mean 0.2955, SD 0.0073, 5th--95th percentile 0.2894--0.3016). This narrowness
is the first indication that an age-count term will behave almost like a
population term.

Candidate-park distance is available for all 2,536 common cells: 537 are within
500~m, 1,390 are within 1,000~m (including the 537), and 1,146 are beyond
1,000~m. The mean distance is
1,141.6~m (median 909.0~m; 95th percentile 2,922.2~m; maximum 5,756.0~m).
The distribution is right-skewed, and the 1,000~m cap used by M4 compresses 45.2\%
of the common sample into the maximum-deficit category. Figure~\ref{fig:environment}
maps LST and NDVI; Figure~\ref{fig:access} maps population and candidate-park
distance. The maps show the expected broad pattern of hotter, less vegetated
central and northern areas, but they also show substantial cell-to-cell
heterogeneity that a city-level summary would obscure.

\begin{figure*}[!t]
  \centering
  \includegraphics[width=\textwidth]{Figures/fig02_environment.pdf}
  \caption{Land surface temperature (a) and NDVI (b) over the 2,536-cell common
  sample. Grey cells are excluded from the common comparison (229 cells). LST is
  the declared 2024 clear-sky morning median, aggregated to 500~m; it is a
  surface temperature, not air temperature or personal exposure. Colour limits
  are the 5th and 95th percentiles of the common sample. Source:
  \texttt{merged\_grid.csv} and \texttt{priority\_models.csv} from the frozen RC1
  artifact; figure generated for this manuscript from those tables.}
  \label{fig:environment}
\end{figure*}

\begin{figure*}[!t]
  \centering
  \includegraphics[width=\textwidth]{Figures/fig03_population_access.pdf}
  \caption{Modelled total population on a logarithmic scale (a) and modelled
  walking-network distance to the nearest OpenStreetMap candidate park, capped at
  1,000~m (b). Population is a 2024 WorldPop estimate; park distance is a
  19 November 2026 candidate proxy. Grey cells are excluded. Distance is a
  modelled route to a representative park point, not a verified entrance or a
  measured public-access trip. Source: \texttt{merged\_grid.csv} and
  \texttt{priority\_models.csv}.}
  \label{fig:access}
\end{figure*}

\subsection{Baseline priority patterns under M1--M4}

Each model assigns 254 cells to its top decile. The score distributions shift
upward and become narrower as components are added: M1 has mean 0.598 (SD 0.268),
M2 0.585 (0.247), M3 0.631 (0.227) and M4 0.669 (0.173). The compression is
partly a consequence of adding components with different distributions and
partly a consequence of percentile clipping. M4 has only two cells tied at the
maximum, compared with 127 for M1, so the top of the M4 ranking is less
saturated.

The geography of the top decile changes substantially across models. Under M1,
the top decile is concentrated in East Jakarta (88 cells, 34.6\% of the top set)
and North Jakarta (77, 30.3\%), with 48 in South Jakarta, 39 in West Jakarta and
only 2 in Central Jakarta. Under M4, East Jakarta still contributes 87 cells
(34.3\%), but West Jakarta rises to 63 (24.8\%), Central Jakarta to 17 (6.7\%),
while South Jakarta falls to 32 (12.6\%) and North Jakarta to 55 (21.7\%). The
M4 top decile contains 1.54 million people, or 14.6\% of the common-sample
population, and 63.4~km$^2$ of land. M1's top decile contains 1.37 million people
and 62.6~km$^2$. The difference in population share is modest, but the
geographic redistribution is not.

Figure~\ref{fig:priority} maps the M1 and M4 top deciles. The two maps share a
broad belt of high priority in the north-east and along parts of the eastern
and southern fringe, but they disagree on many cells in the west and centre.
This is the visual counterpart of the low Jaccard overlap reported below. The
maps are screening maps: a red cell is a candidate for further investigation,
not a validated intervention site.

\begin{figure*}[!t]
  \centering
  \includegraphics[width=\textwidth]{Figures/fig04_priority_maps.pdf}
  \caption{Top-decile priority sets under M1 (LST only; a) and M4 (heat,
  vegetation, age count and park deficit; b). Red cells are the 254 highest-ranked
  cells in each model; light grey cells are eligible but outside the top decile;
  dark grey cells are excluded from the common comparison. The maps use the same
  extent and classification, so the redistribution of priority is directly
  visible. M4 is a provisional screening index, not a validated vulnerability or
  intervention model.}
  \label{fig:priority}
\end{figure*}

\subsection{Ranking displacement and top-priority change}

Model agreement is high in rank correlation but modest in top-set overlap.
Spearman's $\rho$ is 0.936 between M1 and M2, 0.905 between M1 and M3, and 0.847
between M1 and M4. Kendall's $\tau_b$ follows the same order (0.782, 0.736 and
0.658). The M3--M4 correlation is 0.898 ($\tau_b=0.736$), which might suggest
substantial agreement; the top-decile Jaccard is only 0.411. The difference
between these statements is the difference between a global rank correlation and
a boundary-sensitive set comparison. A correlation of 0.9 can coexist with a
large change in which cells cross the top-decile threshold.

The M1-to-M4 mean absolute rank shift is 307 positions (median 239; 90th
percentile 696.5; maximum 1,714) on a 2,536-cell ranking. Almost half of the cells
(48.3\%) move by more than 10\% of the ranking, and only 0.3\% retain the same
rank. The M3-to-M4 shift is smaller but still large: mean 266.7, median 248,
90th percentile 487, maximum 1,065. M1-to-M2 shifts are smaller (mean 199.6,
median 156), and M2-to-M3 shifts are the smallest (mean 129.0, median 88). The
ordering is intuitive: adding vegetation to LST changes the ranking less than
adding population, and adding population changes it less than adding park
distance and age.

Top-set overlap depends strongly on the budget. At the 5\% budget, M1--M4
Jaccard is 0.316, with 61 of 127 cells shared (48.0\% retention); at 10\% the
Jaccard is 0.384, with 141 of 254 cells shared (55.5\% retention); at 20\% the
Jaccard is 0.488, with 333 of 508 cells shared (65.6\% retention). M3--M4
Jaccard is 0.419, 0.411 and 0.532 at the same budgets. M2--M3 is the most
stable pair (0.649, 0.688 and 0.758). The top-decile comparison is therefore not
a generic statement that "the models disagree"; it is a statement that the
disagreement is concentrated near the priority boundary, where small score
changes can move a cell in or out of the set. Jaccard and retention answer
different questions: Jaccard divides the shared cells by the union, whereas
retention divides them by the size of one set. Figure~\ref{fig:rank} shows the displacement and its distribution.

\begin{figure*}[!t]
  \centering
  \includegraphics[width=\textwidth]{Figures/fig05_rank_shift.pdf}
  \caption{Rank displacement between M1 and M4 on the common sample. (a) Each
  point is a grid cell; the dashed line is the identity. Dotted lines mark the
  top-decile boundary (rank 254). (b) Distribution of signed rank change
  (M4 rank minus M1 rank); the median signed change is $-14$ positions, while
  the median absolute change is 239 positions. The high rank correlation
  (Spearman $\rho=0.847$) coexists with a large top-set change (Jaccard 0.384).
  Source: \texttt{priority\_models.csv}.}
  \label{fig:rank}
\end{figure*}

\input{Tables/table5_comparison}

\subsection{Attribution of M3-to-M4 change}

The exact Shapley decomposition attributes 85.46\% of the absolute M3-to-M4 rank
movement to park addition, 13.25\% to reweighting and 1.29\% to the age
substitution. The shares are normalized absolute contributions; they describe
the model change, not the real-world importance of parks, age or vegetation. The
park share is large because the park term is new, enters at a 0.25 weight, and
has a wide 0--1 range that is not compressed by percentile clipping. The age
share is small because the age count is nearly proportional to total population
and the weight change is modest.

The matched-weight incremental comparisons are relative to M3 and are reported
in Table~\ref{tab:attribution_m3}. Adding the park term while holding the other
weights fixed yields a top-decile Jaccard of 0.490 against M3 (167 of 254 cells
shared). Reweighting alone yields 0.961 (249 shared), and the age swap alone
yields 0.984 (252 shared). Adding park plus reweighting yields 0.415; age plus
park yields 0.499; all three changes yield 0.411. These values answer how much
each change moves the M3 top set; they are not similarities to M4.

The same scenarios compared with M4 are reported in
Table~\ref{tab:attribution_m4}. Park addition alone is the closest
single change to M4 (Jaccard 0.861; 235 of 254 cells shared), while reweighting
alone (0.423; 151 shared) and the age swap alone (0.419; 150 shared) are much
further away. Park plus reweighting reaches 0.977, and age plus park reaches
0.847. Read together, the two tables show that the park term is the dominant
source of set change, but the numbers are not interchangeable: Jaccard divides
shared cells by the union of the two sets, whereas retention divides them by one
set's size.

The construct ablations relative to M4 are also in Table~\ref{tab:attribution_m4}.
A park-only index has a top-decile Jaccard of 0.028 with M4 (14 of 254 cells
shared), showing that M4 is not merely a park map; a no-park variant has Jaccard
0.423, and a heat-only variant 0.384. An equal-weight variant has Jaccard 0.868,
and a proportional M3-plus-park variant 0.924. These benchmarks are internal
comparisons on the same sample, not external validation.
Figure~\ref{fig:attribution} shows the Shapley shares and the scenario-family
envelope.

\begin{figure*}[!t]
  \centering
  \includegraphics[width=\textwidth]{Figures/fig06_attribution_sensitivity.pdf}
  \caption{(a) Exact three-factor Shapley shares of absolute M3-to-M4 rank
  movement: park addition 85.46\%, reweighting 13.25\%, age substitution 1.29\%.
  (b) Jaccard overlap of the M4 top decile with each scenario family at the 10\%
  budget; horizontal lines span the minimum--maximum range. The joint factorial
  and route-policy families contain the lowest values. These are model-change
  decompositions, not causal effects. Source:
  \texttt{rank\_change\_attribution.csv} and
  \texttt{robustness\_family\_summary.csv}.}
  \label{fig:attribution}
\end{figure*}

\input{Tables/table6a_attribution_m3}
\input{Tables/table6b_attribution_m4}

\subsection{Accessibility-assumption sensitivity}

Park-access assumptions change the top decile materially. The positive-tag-only
destination scenario, which retains only 10 destinations, produces a top-decile
Jaccard of 0.420 against the baseline when eligibility changes and 0.425 when
anchors are fixed. This is a stress test, not a better estimate: the positive-tag
set is sparse, and 2,482 cells are saturated at the 1,000~m cap. The conservative
connector scenario, which removes 24,370 directed edges and increases missing
routes to 407, produces a fixed-set Jaccard of 0.896. The baseline itself has 226
missing routes. The difference between 0.420 and 0.896 is the difference between
a destination-coverage stress and a permission stress; both are plausible under
different evidence assumptions, and neither is ground truth.

The missing-route bounds are also informative. Assigning the 1,000~m cap to
missing routes (the worst-case deficit) and assigning the observed median
distance (a neutral alternative) produce different top sets; the scenario family
spans Jaccard 0.906 at the 10\% budget. This is smaller than the destination
stress but not negligible. The 13 route-missing cells that could enter the top
254 under a zero-deficit assumption are a further boundary uncertainty; they are
not included in the common set and are mapped in the supplementary material.
The 18 cells that are in the full-thermal M1 top 277 but lost from the common set
are all route-missing (15 North, 2 East, 1 South). These cells are not evidence
that the common set is representative; they are evidence that route missingness
is spatially structured.

\subsection{Demographic-specification sensitivity}

The demographic specification matters more for interpretation than for the
primary ranking. Replacing the age count with total population at the same weight
leaves the top decile almost unchanged (Jaccard 0.977; 251 of 254 retained). This
is consistent with the near-collinearity of the two counts (Spearman
$\rho=0.9995$). The age count is therefore a population-weighted exposure term,
not a distinct age-vulnerability signal. Replacing the count with a vulnerable
share changes the set substantially (Jaccard 0.347; 131 retained). Child share
yields 0.451 and elderly share 0.392. The share variants answer a different
question: where vulnerable groups are a larger fraction of residents, not where
many vulnerable people live. The primary M4 answers the latter question only.

The distinction is visible in the data. The vulnerable share is almost constant
across the common sample (mean 0.2955, SD 0.0073), and its correlation with LST
is essentially zero (Spearman $\rho=0.0057$). The top-M4 cells have a vulnerable
share of 0.2894--0.3016, a range of 1.2 percentage points. A model that claimed to
identify age-vulnerable heat inequality on the basis of this term would be
overstating what the data contain. The share scenario is useful precisely because
it shows that a composition-oriented objective produces a different geography;
it is not a correction that makes M4 better.

\subsection{The mechanical park-threshold effect}

Under the provisional M4, the top-decile cutoff is 0.873739849. The maximum M4
score among the 537 common cells within 500~m of a candidate park is 0.869648776.
The cutoff exceeds that maximum by 0.004091, so no near-park cell enters the top
254. The closest top-decile cell is 638.2~m from a candidate park; 174 of the 254
top cells are at the 1,000~m cap, and the remaining 80 are between 638 and
1,000~m. This is a property of the additive index and the observed input
distribution, not an independently validated empirical discovery that the most
heat-vulnerable places lack parks.

A concrete case illustrates the arithmetic. Cell \texttt{JKT\_R056\_C048} in
North Jakarta has an LST of 45.58~$^\circ$C, a modelled population of 6,971 and a
candidate-park distance of 487.96~m. Its M4 score is 0.869649 and its rank is 274,
just outside the top 254. Of the 254 top cells, 236 are cooler than this cell.
The model is not "wrong" in the sense of an arithmetic error; it is doing what
the additive formula specifies. But the result is a screening priority that
penalizes a hot, populous, near-park cell because the park term contributes only
0.122 while the other components do not compensate enough to cross the cutoff.
Presenting this as evidence that parks do not matter, or that near-park
residents are not heat-vulnerable, would invert the meaning of the model.

\subsection{Missing-data exclusions and representativeness}

The 229 excluded cells account for 8.28\% of the master grid and 4.22\% of the
modelled population (463,697 people). They are not a random subset. Included
cells have a mean LST of 41.486~$^\circ$C, mean NDVI 0.315, mean coverage 0.940
and mean snapping distance 49.98~m. Excluded cells have a mean LST of
40.283~$^\circ$C, mean NDVI 0.365, mean coverage 0.821 and mean snapping distance
98.81~m. The excluded cells are therefore cooler, greener, more peripheral and
less well connected (Table~\ref{tab:selection}). Their city distribution is East Jakarta 96, North Jakarta
77, South Jakarta 28, West Jakarta 22 and Central Jakarta 6. Most are
route-missing (226 of 229), with two population-only exclusions and one low-LST
exclusion.

This profile has two consequences. First, the common-set estimate is conditional
on the included cells; it is not a Jakarta-wide estimate. Second, the excluded
cells are not uniformly low-priority: 18 cells that are in the full-thermal M1
top 277 are lost from the common set, and 13 route-missing cells could enter the
top 254 under a zero-deficit assumption. The common-set map should therefore be
read alongside the exclusion map in the supplementary material. A priority map
that omits the excluded cells without showing them invites the reader to mistake
the study area for the city.

\input{Tables/table9_selection}

\subsection{Global and local spatial clustering}

Global Moran's $I$ is positive and highly significant for all four models under
the permutation convention: M1 0.6711, M2 0.6540, M3 0.6644 and M4 0.5930, with
$p_{\mathrm{sim}}=0.001$ in every case. The expected value under randomization is
$-0.000395$. The M4 value is lower than M1--M3, consistent with the addition of a
park term that is less spatially smooth than LST. The $p$-value of 0.001 is the
minimum attainable at 999 permutations; it should be read as $p<0.002$, not as
$p=0$.

Local cluster labels are materially convention-dependent. Under the provisional
directed/999/BH convention, the M4 map has 319 high--high, 323 low--low, 11
high--low and 10 low--high cells, with 663 significant cells and 1,872
non-significant cells. Two-sided testing at 999 permutations yields 496
significant cells, a change of 167 labels. Increasing permutations to 9,999
yields 687 significant cells under directed testing and 493 under two-sided
testing. BY correction at 999 permutations yields zero significant cells, because
the minimum attainable $p$-value (0.001) exceeds the BY threshold
($0.05/2{,}535\approx1.97\times10^{-5}$); at 9,999 permutations it yields 284
directed and 234 two-sided significant cells. The zero is a resolution artifact,
not evidence that no local clustering exists.

Figure~\ref{fig:lisa} maps the M4 clusters under the provisional convention and
shows the count range across conventions. The map is useful for exploring
spatial structure, but it is not a validated cluster map. A different tail,
permutation count or correction changes which cells are labelled. We therefore
report the full range and avoid statements such as "significant everywhere" or
"no significant clusters" without naming the convention.

\begin{figure*}[!t]
  \centering
  \includegraphics[width=\textwidth]{Figures/fig07_lisa.pdf}
  \caption{Local Moran's $I$ for M4. (a) Cluster labels under the provisional
  directed/999-permutation/Benjamini--Hochberg convention. (b) Number of
  significant local clusters under alternative tail, permutation and correction
  conventions. The zero under BY/999 is a minimum-attainable-$p$ artifact, not
  evidence of no clustering. Local labels are not causal or health findings.
  Source: \texttt{lisa\_m4.csv} and
  \texttt{lisa\_tail\_permutation\_sensitivity.csv}.}
  \label{fig:lisa}
\end{figure*}

\input{Tables/table8_lisa}

\subsection{Scenario envelope and stability}

The scenario envelope shows partial stability, not uniform robustness. At the
10\% budget, the aggregation family has Jaccard 0.881; QA exclusions 0.969--0.996;
missing support 1.000; thermal support 1.000; land cover 0.782; normalisation
0.734--0.814; park threshold 0.521--0.682; weights 0.868--0.954; population
definition 0.347--0.977; route policy 0.420--1.000; and the 16-scenario joint
factorial 0.583--0.728. The joint factorial never reaches 0.8, and the route and
population families contain the lowest values. A claim that "the priorities are
robust" is therefore not supported. A more accurate statement is that a subset
of cells is stable across several families, while the boundary of the top decile
is contingent on construct, weight, demographic and access choices.

The stability map in the supplementary material shows this subset. It should not
be read as a validated core: stability across the scenarios we happened to run
is not stability across all plausible models. The scenario registry itself is a
choice, and two families (entrance/barrier and thermal sampling) are blocked by
missing evidence. The envelope is an honest description of the sensitivity we
could test, not a complete uncertainty budget.

\input{Tables/table7_sensitivity}

\subsection{Summary of the evidence}

Three results are robust in the sense that they follow from exact arithmetic and
do not depend on interpretive choices. First, M1 and M4 produce materially
different top-decile sets (Jaccard 0.384) despite a high rank correlation
(0.847). Second, the age count is nearly collinear with total population
($\rho=0.9995$), and a total-count swap leaves the top decile almost unchanged
(0.977), while a share swap changes it substantially (0.347). Third, under the
observed M4 inputs, no cell within 500~m of a candidate park enters the top 254;
the cutoff and the maximum near-park score differ by 0.0041.

Three further results are conditional. First, the attribution shares (85.5\%
park, 13.2\% reweighting, 1.3\% age) describe this M3-to-M4 configuration and
should not be generalised to other weights or caps. Second, the access stress
tests show that destination and permission assumptions can change the top set
materially, but the positive-tag scenario is not ground truth. Third, the LISA
labels are convention-dependent, and the range (0 to 687 significant cells) is
too wide to support a single cluster narrative. These conditional results are
still useful: they tell the analyst where the index is fragile and what evidence
would be needed to resolve the fragility.

\section{Discussion}

\subsection{Adding indicators changes the question, not just the map}

The four models are often described as progressively "better" because they
include more information. That framing is misleading. M1 answers a narrow
question: where is the surface hottest? M2 asks where heat coincides with low
vegetation. M3 asks where heat and vegetation coincide with many residents. M4
asks where heat, vegetation, a count of children and older adults, and a modelled
park deficit combine into a high score. These are different questions, and none
is a gold standard for the others. A cell that ranks first under M1 may rank
274th under M4 without either model making an arithmetic error. The appropriate
question is not "which model is right?" but "which question is the decision
supposed to answer, and how transparent is the value judgement embedded in the
composite?"

This is especially important for policy communication. A heat-only map can be
read as a physical hazard map. A composite map is read by many audiences as a
vulnerability or priority map, even when its authors intend it as a screening
device. The label matters. We therefore describe M4 as a provisional
heat--population--park-deficit screening index, not as a heat inequality or
age-vulnerability index. The distinction is not semantic: it determines which
claims are licensed by the output.

\subsection{Why the park term dominates}

The park term dominates the M3-to-M4 rank movement for three structural reasons.
First, it is new: the other components are present in M3 in some form, so the
park term introduces a dimension that was previously absent. Second, it enters at
a 0.25 weight, which is large relative to the 0.30/0.20 weights on LST and
vegetation. Third, it is defined on a fixed 0--1 scale with a hard cap, whereas
the other components are percentile-normalised and compressed near their central
range. A cell at the cap receives the full 0.25 contribution; a cell at 500~m
receives 0.125. The additive form makes these contributions tradeable against
heat and population. The result is a ranking in which distance to a candidate
park can outweigh a large LST difference.

None of this demonstrates that parks are unimportant. It demonstrates that the
index is sensitive to the park term's scale and weight. A different cap, a
nonlinear distance decay, a binary threshold or a separate park-access model
would produce a different ranking. The park term is a legitimate planning
construct, but its dominance is a design property that should be disclosed, not
a finding about the world. The near-park hot cell example is the clearest
illustration: the model ranks a 45.6~$^\circ$C, 6,971-person cell outside the top
254 because its park deficit is small. A policymaker who reads the top decile as
"the most heat-vulnerable places" would draw the wrong conclusion. A
policymaker who reads it as "places where a particular additive screening rule
scores high" would not.

\subsection{What the age term does and does not contribute}

The age count contributes little beyond total population in this dataset. The
two counts have a Spearman correlation of 0.9995, and replacing one with the
other changes only 3 of 254 top cells. This is not because age is unimportant;
it is because the age count is a population-scaled measure in a setting where the
vulnerable share is nearly constant. A count answers "where are many vulnerable
people?" and will therefore track population. A share answers "where are
vulnerable people a larger fraction of residents?" and produces a different
geography. The share variants have Jaccard 0.347--0.451 with M4, showing that the
choice of estimand matters more than the choice of age group.

The policy implication is to separate the questions. If the goal is to find
places with many older adults during a heat event, a count is defensible, but it
should be labelled as population-weighted exposure rather than age vulnerability.
If the goal is to test whether age composition is associated with heat, a share
or a regression with population controls is needed. Mixing a count and a share in
one composite, or interpreting a count as a composition, obscures the
distinction. Our results support the first interpretation for M4 and reject the
second.

\subsection{Why park-access evidence is uncertain}

The park layer is a candidate inventory, not a verified public-access inventory.
No destination has verified public access or a verified entrance; 980 of 1,015
are tagged unknown, and all 1,015 remain unresolved. OSM completeness and
tagging quality vary, and a missing tag is not evidence of a closed park
\citep{haklay2010,barron2014,herfort2023}. The walking graph is also a model:
origins and destinations are snapped to nearest nodes, park polygons are
represented by representative points, and 167 origin connectors plus 78 park
connectors conflict with explicit node restrictions. The conservative scenario
removes 24,370 directed edges and increases missing routes to 407. The
positive-tag-only scenario retains 10 destinations and saturates 2,482 cells at
the cap. Neither scenario is a validated access measure.

This uncertainty is not a reason to discard the park term. It is a reason to
report it as a proxy and to show the range. The baseline and conservative
scenarios agree on a large share of the top decile (Jaccard 0.896), which is
reassuring for the permission dimension. The destination stress is much larger
(Jaccard 0.420), which is a warning that the candidate inventory itself is a
major source of uncertainty. A field audit of access and entrances for the
highest-impact destinations would reduce this uncertainty more than a more
complex routing algorithm on unverified tags.

\subsection{Sample selection and geographic generalisation}

The common sample excludes 229 cells (8.3\%) and 4.2\% of the modelled
population. The excluded cells are cooler, greener, more peripheral and less
well connected. This is not a random loss: it is a spatially structured
selection driven by route missingness and source support. The common-set
estimate is therefore conditional on the included cells, and the excluded cells
should be shown on every priority map. The 18 full-thermal M1 top cells lost from
the common set and the 13 route-missing cells that could enter the top 254 are
not footnotes; they are part of the uncertainty of the map.

The appropriate language is "the five mainland cities under the common-set
configuration", not "Jakarta". A Jakarta-wide claim would require either
complete route coverage or an explicit imputation model for missing routes, and
neither is available. The current artifact supports a conditional screening
statement. It does not support a citywide prevalence estimate, a comparison of
city-level rates, or a claim that the excluded periphery is unimportant. The
periphery is cooler and greener on average, but it also contains hot cells and
people, and the average masks that heterogeneity.

\subsection{Spatial inference and policy interpretation}

The LISA results illustrate why spatial inference needs conventions. The same M4
map yields 663 significant local clusters under directed/999/BH, 496 under
two-sided/999/BH, 687 under directed/9,999/BH and zero under BY/999. The zero is
not a substantive finding; it is a consequence of the minimum attainable
permutation $p$-value and the BY threshold. A paper that reported only the
largest number would overstate certainty; a paper that reported only the zero
would understate it. The honest report is the range, with the primary convention
pre-specified and the alternatives disclosed.

Local clusters are also not causal processes. A high--high cluster of M4 scores
can arise from the index's own spatial structure, from the distribution of
population, or from the park term. It does not identify a social mechanism, a
health effect or a policy lever. The cluster map is a exploratory tool for
generating hypotheses and for designing field work. It is not a substitute for
an outcome model.

\subsection{Comparison with prior work}

Our results overlap with prior work in several ways. The LST--NDVI relationship
and the surface urban heat island literature justify the M1/M2 layers
\citep{oke1982,voogt2003,weng2004,estoque2017}. Heat vulnerability indices
justify the demographic layers \citep{reid2009,wolf2013}. Network-access studies
justify the park layer \citep{oh2007,kim2026,feng2019}. Composite-indicator
research justifies the sensitivity design \citep{saisana2005,paruolo2013,greco2019}.
Jakarta studies establish the local relevance \citep{siswanto2023,putra2021,rizki2024,varquez2020}.
The contribution is not the combination of these ingredients. It is the
quantitative attribution of how the combination changes priorities, the
demonstration that an age count is nearly redundant with population, and the
disclosure that an additive park term can mechanically determine the top decile.

Compared with Sidiqui et al.\ \citep{sidiqui2022}, we use a finer grid, a
network-based candidate-park measure and an explicit attribution decomposition,
but we do not claim a better vulnerability map. Compared with Kim and Kwak
\citep{kim2026}, we do not implement a cooling-weighted G2SFCA or a
demand-allocation model; our park term is a simpler distance deficit, and our
contribution is the sensitivity of the resulting ranking rather than the
sophistication of the access measure. Compared with Fung et al.\
\citep{fung2024}, we do not evaluate heat-mitigation interventions; we evaluate
how a screening index responds to specification choices. These differences are
substantive, and they define a narrower but defensible contribution.

\subsection{What a policymaker may legitimately infer}

A policymaker may legitimately use the maps to generate a shortlist of cells for
field inspection, to compare the spatial implications of different screening
questions, and to see where the ranking is sensitive to assumptions. The maps
can support a conversation about whether a park-access term should be included,
how much weight it should receive, and what evidence would be needed to verify
the candidate destinations. They cannot support a claim that a particular cell
is the best place to invest, that a park would cool it by a specified amount, or
that the top decile corresponds to excess mortality or morbidity. Those claims
require outcome data, cooling-effectiveness estimates, cost data and a causal
design that this study does not provide.

The most defensible policy use is procedural: require a sensitivity report
alongside any composite priority map, show excluded and unknown cells, state the
access construct, and separate screening from evaluation. The artifact's
claim-evidence register is designed for that purpose. It does not replace
scientific judgement, but it makes the judgement inspectable.

\subsection{What remains unknowable from the present data}

Several questions cannot be answered with the current evidence. We do not know
the original 2024 acquisition dates, the seasonal representativeness of the
thermal composite, or whether the 2026 park and network snapshot resembles the
2024 network. We do not know which candidate parks are publicly accessible, when
they are open, whether they have entrances on the modelled network, or whether
they provide cooling refuge. We do not know how many people use them, how far
they actually walk, or how heat affects their health. We do not know whether the
excluded cells would change the top decile if routes were resolved. These are
not minor caveats; they define the boundary between a screening exercise and an
impact evaluation.

\section{Limitations}

The limitations of this study are specific and consequential.

\begin{enumerate}
\item \textbf{Thermal acquisition provenance.} The original Earth Engine scene,
task and AOI records are unavailable. The 2024 label is a declared recipe period,
not a verified acquisition history. Scene counts, dates and seasonal sampling
cannot be reconstructed from the delivered median raster. No annual, seasonal or
trend claim is supported.
\item \textbf{LST is not air temperature or personal exposure.} LST measures the
radiant surface temperature of the skin of the city. It is not the air
temperature a person experiences, and it does not incorporate humidity, wind,
shade, clothing, activity or access to cooling. The relationship between LST
and heat stress is context-dependent \citep{schwarz2012}.
\item \textbf{Temporal composite representativeness.} The median composite is
built from unequal clear-observation counts across pixels (grid mean 5.96--27.4
observations per pixel in the delivered data). The median is not an annual mean,
and the monthly/seasonal balance is unknown.
\item \textbf{Mixed reference periods.} LST and demographic inputs refer to a
declared 2024 period; the OSM park and pedestrian network is a 19 November 2026
snapshot. The analysis is a mixed-period candidate proxy, not a contemporaneous
2024 reconstruction.
\item \textbf{OSM park inventory coverage.} The inventory contains 1,015
candidate polygons, but OSM completeness and tagging vary. Some public parks may
be missing; some mapped polygons may not be parks. The candidate set is not a
census.
\item \textbf{Unverified public access.} No candidate has verified public access
or a verified entrance. The 980 unknown, 28 conditional and 7 positive-tag
classifications are evidence states, not access determinations. A positive tag
is an unverified assertion.
\item \textbf{Unverified park entrances and connectors.} Park polygons are
represented by representative points, and origins and destinations are connected
to the graph at nearest nodes. These connectors are not field-verified. The
conservative scenario removes 24,370 directed edges and leaves 407 missing
routes.
\item \textbf{Walking-network approximations.} The graph is a pedestrian network
derived from OSM. It does not encode opening hours, temporary closures, stairs,
crossings, security gates, flooding or construction. Shortest-path distance is
not the same as walkability or perceived access.
\item \textbf{Missing park distances.} The baseline has 226 route-missing cells
(8.2\%). They are excluded, not assigned a distance. The conservative scenario
has 407. The 13 route-missing cells that could enter the top 254 are a lower
bound on the effect of missingness.
\item \textbf{Demographic raster support.} Six cells remain unsupported in total
population and both age groups, including two substantively large North Jakarta
cells. They are retained as nulls, not zeros. The 13 recovered cells depend on
buffered age surfaces and a hybrid aggregation rule.
\item \textbf{Population aggregation.} The primary hybrid rule is not additive
under repartitioning; the fractional alternative changes 16 top-M4 cells and has
Jaccard 0.881 with the baseline. The choice of primary allocation is unsigned.
\item \textbf{Partial coastal grids.} 383 cells are partial polygons; 65 have
less than 10\% coverage. Small-area population and density estimates are
uncertain, and the lattice clipping affects the spatial graph.
\item \textbf{Weighting and normalisation.} The weights are provisional and not
derived from Jakarta-specific evidence. Percentile clipping saturates scores and
creates ties. The park cap is a modelling choice. Different weights, caps or
normalisation rules produce different rankings, as the scenario envelope shows.
\item \textbf{Sample selection.} The common sample excludes 229 cells (8.3\%)
and 4.2\% of the modelled population. Excluded cells are cooler, greener, more
peripheral and less connected. The common-set estimate is conditional, not
Jakarta-wide.
\item \textbf{Spatial inference conventions.} LISA labels depend on tail,
permutation count and multiple-testing correction. The significant count ranges
from 0 to 687 across the conventions we tested. The 0.001 permutation $p$-value
is a resolution limit, not $p=0$.
\item \textbf{Limits of causal interpretation.} The models are spatial
associations. They do not identify causal effects of parks, vegetation,
population or age on heat exposure or health. They do not estimate intervention
effectiveness, cost-effectiveness or distributional welfare.
\item \textbf{Raw-data portability.} The raw WorldPop/OSM/GraphML bundle is not
tracked in Git; exact raw-input hashes require the separate frozen bundle. RC1's
original clean-checkout test failure and three Windows line-ending receipt
mismatches were remediated in the RC2 engineering candidate: 236 tests plus 7
subtests pass in a clean checkout with zero failures/errors/skips; RC1 and RC2
receipts validate 9/9 under the canonical Git-blob SHA-256 policy; and the
standard and remediation output manifests validate 69/69 and 80/80 relative
paths. RC2 is an engineering remediation, not scientific approval, and it does
not make the raw bundle public.
\end{enumerate}

\section{Conclusion}

This paper examined how spatial priorities for urban heat screening in Jakarta
change when LST and vegetation are supplemented with population, age count and
modelled park access. The answer is: substantially, and in ways that are not
captured by a single correlation or a single map. M1 and M4 share 141 of their
254 top-decile cells (55.5\% retention), but the Jaccard overlap is only 0.384
because each set contributes 113 unique cells; the mean absolute rank shift is
307 positions. The
M3-to-M4 change is dominated by the park term (85.5\% of absolute rank movement),
with reweighting (13.2\%) and the age substitution (1.3\%) playing smaller roles.
The age count is nearly collinear with total population; a count-to-count swap
leaves the top decile almost unchanged, while a count-to-share swap changes it
substantially. Under the observed M4 inputs, no cell within 500~m of a candidate
park enters the top 254, because the top-decile cutoff exceeds the maximum
near-park score. That is a property of the additive index, not evidence that
parks are unimportant.

The study's contribution is therefore an attribution and uncertainty audit, not
a new integrated framework and not a validated vulnerability map. It shows which
conclusions are robust to the assumptions we could test, which are contingent,
and which cannot be evaluated with the available evidence. The most important
negative findings are that an age-count term does not measure age composition,
that an additive park-deficit term can mechanically determine the top decile,
and that local cluster labels are convention-sensitive. These findings are
useful precisely because they constrain what a screening map can be used to say.

The next scientific steps are clear. Recover or formally restrict the thermal
acquisition lineage. Field-verify public access and entrances for the
highest-impact candidate parks. Resolve or explicitly model the 226--407 missing
routes. Separate the heat-plus-population question from the age-composition and
park-access questions, rather than forcing one composite to answer all three.
Pre-specify the LISA convention and report the range. Until those steps are
completed, the present results should be read as a reproducible exploratory
screening artifact and a working preprint draft, not as an approved scientific
release or a policy recommendation.

\section*{Data and Code Availability}

The analysis is built on the frozen RC1 artifact of the Jakarta Heat Inequality
Mapper, commit \texttt{\seqsplit{3b8ed05be92c35bbb2b3a617688540ee856b11c8}}. The artifact
contains the master grid, processed remote-sensing and demographic tables, the
priority models, the scenario tables, the LISA outputs and the reproduction
scripts. The canonical manuscript source, bibliography, figure-generation
scripts, claim-evidence register and build instructions are in the
\texttt{paper/zenodo\_preprint} directory of the project repository.

The repository URL and the public archive DOI are \textbf{not yet assigned}.
A Zenodo deposit for this manuscript has not been created, and no DOI is
claimed. The RC1 artifact is an internal, conditionally released research
candidate; it is not a public data release. The WorldPop input is subject to its
own CC BY 4.0 terms \citep{worldpop2025}; OpenStreetMap data are subject to the
Open Database License and attribution requirements \citep{osm2026}. Users must
verify the source licences before redistribution.

The exact reproduction environment is Python 3.12.10 on Windows x64 with pandas
3.1.1, numpy 2.5.5, libpysal 4.15.1, esda 2.10.0, geopandas 1.2.0 and scipy
1.18.1. The RC1 reproduction commands are documented in
\texttt{docs/RC1\_REPRODUCE.md}. RC2 engineering evidence is documented in
\texttt{rc2\_evidence/RC2\_REPRODUCE.md}: source revision
\texttt{ac656e5b1b123d3bb5c5f618489682109bba2de9}, artifact revision
\texttt{267ebf78efc985926abd2ad64ec4d669138e1a6e}, implementation SHA-256
\texttt{1858264d4710f0209cd8fa03320af541ec99ce10f1bc31635aac5a24f826207a},
236 tests plus 7 subtests passing in a clean checkout with zero
failures/errors/skips, 9/9 RC1 and RC2 receipts under the canonical Git-blob
SHA-256 policy, 69/69 standard and 80/80 remediation output-manifest entries
with repository-relative paths, and a 63-table replay with maximum numeric
difference 0.0. RC2 leaves the model definitions and empirical outputs identical
to RC1 and remains scientifically unsigned; the raw WorldPop/OSM/GraphML bundle
is still distributed separately and is not tracked in Git. The
manuscript's PDF build uses the Manuscript Multi-Target LaTeX Template v1.0.0
\citep{aramayo2026}; the template is GPL-3.0-or-later and is attributed in the
repository's \texttt{SOURCE\_PROVENANCE.md}.

\section*{Author Contributions}

Author contributions are \textbf{to be completed and approved by the three named
authors}. Suggested CRediT categories for the final version include:
conceptualization; data curation; formal analysis; investigation; methodology;
software; validation; visualization; writing---original draft; writing---review
and editing. No AI system is listed as an author. The human authors must confirm
all contributions before submission.

\section*{Funding}

No funding source has been verified for this manuscript. Funding and
acknowledgement statements are \textbf{to be completed} by the project
coordinator. Do not insert a grant number unless it has been verified against an
award record.

\section*{Conflicts of Interest}

No conflicts of interest have been declared. This statement is a placeholder
pending written confirmation from all authors.

\section*{Acknowledgements}

The manuscript uses the Manuscript Multi-Target LaTeX Template v1.0.0 by Rodolfo
Aramayo (Zenodo DOI \texttt{10.5281/zenodo.22018962}), distributed under
GPL-3.0-or-later. The template's class and bibliography style derive from
Ricardo Henriques's HenriquesLab bioRxiv template; upstream attribution is
retained in the template licence and provenance files. The template is a
third-party typesetting resource; its use does not imply Zenodo endorsement or
acceptance. We thank the open-source contributors to Landsat, WorldPop,
OpenStreetMap, Python, GeoPandas, libpysal and esda. This acknowledgement does
not imply that any named individual has reviewed or approved the manuscript.

\section*{Supplementary Material}

The supplementary PDF contains: (S1) a data-quality and missingness register;
(S2) scenario-family results at 5\%, 10\% and 20\% budgets; (S3) city-level
composition of the top deciles; (S4) alternative LISA summaries; (S5) the
exclusion map and score distributions; (S6) all 49 executed scenario
comparisons; (S7) the six demographic no-support cells; (S8) the
master-to-common eligibility reconciliation; and (S9) figure and table
provenance with the machine-readable supplementary files. The main text reports
the results needed for the argument; the supplement documents the sensitivity
detail and source support that would otherwise interrupt it.

\clearpage
\section*{References}
\markboth{References}{References}
\begingroup
\setlength{\emergencystretch}{5em}
\sloppy
\bibliography{References}
\endgroup

\end{document}



